\chapter*{Introduction Générale}
\addcontentsline{toc}{chapter}{Introduction Générale}

\section*{Contexte}
La gestion de la scolarité d'un établissement d'enseignement repose sur un ensemble d'opérations
quotidiennes : inscription et mise à jour des dossiers des élèves, constitution des sections et des
groupes, élaboration des emplois du temps, saisie des présences, consignation des observations et des
sanctions, organisation des examens et publication des notes. Ces opérations produisent chacune un
document, et ces documents sont traditionnellement tenus sur papier ou répartis dans des classeurs
tableurs propres à chaque service. La conséquence est connue : la même information existe en plusieurs
exemplaires qui divergent, une recherche dans les archives prend plusieurs minutes, et le responsable
pédagogique ne dispose jamais d'une vue d'ensemble au moment où il doit décider.

La transformation numérique du secteur éducatif apporte une réponse à cette fragmentation, mais elle se
heurte dans de nombreux établissements à une contrainte matérielle : l'absence de connexion Internet,
ou l'interdiction d'exposer des données nominatives d'élèves sur un réseau public. Une application
hébergée sur un serveur web distant est alors inutilisable, et une application purement locale sur un
poste unique ne permet pas le travail simultané de plusieurs services.

\textbf{\appname} a été conçue pour ce contexte précis. Il s'agit d'une application de bureau qui
s'installe sur chaque poste de travail et se connecte, à travers le réseau local de l'établissement, à
une base de données PostgreSQL unique installée sur un serveur interne. Aucune requête ne quitte le
réseau : ni les polices de caractères, ni les rapports imprimés, ni les photos des élèves ne dépendent
d'un service externe. Le poste de travail ne fait que présenter et saisir l'information ; la vérité est
centralisée, partagée et sauvegardée sur le serveur.

\section*{Problématique}
La gestion actuelle, lorsqu'elle s'appuie sur le papier et sur des fichiers dispersés, présente quatre
limites structurantes. D'abord, la \textbf{duplication} : un élève change de section, mais son ancienne
affectation reste inscrite dans le registre du secrétariat et dans le tableau du professeur. Ensuite,
la \textbf{lenteur} : retrouver l'historique des absences d'un élève avant un conseil de discipline
suppose de feuilleter plusieurs registres. Puis l'\textbf{absence de contrôle} : rien n'empêche un
utilisateur non habilité de modifier une note ou de supprimer une sanction, et aucune trace ne
conserve l'identité de celui qui a écrit une ligne. Enfin, la \textbf{fragilité} : un classeur égaré ou
un disque défaillant efface des années de scolarité, faute de sauvegarde organisée.

\textbf{Question de recherche : comment concevoir et développer une application de gestion académique
fonctionnant intégralement hors ligne sur le réseau local d'un établissement, qui centralise les
dossiers des élèves, la structure pédagogique, les emplois du temps, les évaluations et la vie
scolaire, tout en garantissant la sécurité des accès et la traçabilité des opérations ?}

\section*{Objectifs de l'Application}
L'objectif principal de cette application, nommée \textbf{\appname}, est de fournir aux services d'un
établissement un outil unique de gestion de la scolarité, utilisable sans connexion Internet, dont
chaque profil ne voit et ne modifie que ce qui relève de sa responsabilité.

Les objectifs spécifiques sont les suivants :
\begin{itemize}
    \item Remplacer les registres papier par un dossier numérique unique par élève, couvrant l'identité,
          l'état civil, la santé, la famille, les contacts d'urgence, la scolarité et les unités
          d'enseignement ;
    \item Structurer l'établissement dans la base : année d'études, semestres, semaines, filières,
          sections, groupes, salles et créneaux, afin que l'emploi du temps et les examens s'appuient sur
          un référentiel unique ;
    \item Enregistrer la vie scolaire au fil des séances : présences, justifications d'absence, motifs
          d'observation, récompenses et sanctions, avec leur historique consultable ;
    \item Organiser les examens, la répartition des candidats, la surveillance et la saisie des notes,
          puis calculer les moyennes semestrielles selon les coefficients ;
    \item Constituer une bibliothèque documentaire partagée : livres, romans, leçons et vidéos,
          avec catégories, publics, revue éditoriale et notifications ;
    \item Centraliser les informations dans une base de données unique, sauvegardée sur le serveur de
          l'établissement ;
    \item Garantir la sécurité des données grâce à un système d'authentification, à une gestion des
          comptes par l'administrateur et à un contrôle d'accès par rôle vérifié côté serveur ;
    \item Offrir une interface ergonomique, intuitive et accessible, disponible en anglais, en arabe et
          en français, avec un affichage adapté à la lecture de droite à gauche.
\end{itemize}

\section*{Organisation du Rapport}
Le rapport de ce projet est organisé de la façon suivante :
\begin{itemize}
    \item \textbf{Chapitre 1}, intitulé \textbf{Conception et Présentation de l'Application}, est dédié à
          la présentation générale de l'application, au modèle physique de données, au diagramme de cas
          d'utilisation et à l'architecture fonctionnelle basée sur les rôles ;
    \item \textbf{Chapitre 2}, intitulé \textbf{Réalisation de l'Application}, présente l'environnement de
          développement et l'architecture de la solution : processus principal, pont de communication,
          interface, internationalisation, génération des rapports et type d'application ;
    \item \textbf{Chapitre 3}, intitulé \textbf{Guide d'Utilisation de l'Application}, regroupe le guide
          d'installation sur le serveur et sur les postes clients, les fonctionnalités communes, puis le
          guide d'utilisation structuré selon les sept rôles de l'application.
\end{itemize}
